% !TeX program = xelatex
\documentclass[11pt]{article}
\newcommand{\notesdir}{../../lecture-notes}
\providecommand{\notesdir}{.}

% -----------------------------------------------------------------------------
% Shared preamble for the CO5059 lecture notes and the quiz-review documents.
% Documents outside this folder define \notesdir (path to this folder) first.
% -----------------------------------------------------------------------------

% Page layout and typography
\usepackage[a4paper,portrait,margin=1in]{geometry}
\usepackage{fontspec}
\newfontfamily\arialfont{Arial}
\usepackage{microtype}
\usepackage{parskip}

% Mathematics and diagrams
\usepackage{amsmath,amssymb,amsthm,mathtools}
\usepackage{tikz-cd}
\usetikzlibrary{arrows.meta,positioning,shapes.geometric,shapes.arrows,calc,fit,backgrounds}

% Reusable styles for flowcharts and concept diagrams.
\tikzset{
  concept/.style={
    draw=brandprimary,
    fill=brandsecondary!8,
    rounded corners,
    minimum width=3.2cm,
    minimum height=0.9cm,
    align=center,
    font=\small,
  },
  diagram arrow/.style={-{Latex[length=2.5mm]},thick,brandprimary},
}

% Tables, figures, code, and lists
\usepackage{graphicx}
\graphicspath{{\notesdir/}}
\usepackage{enumitem}
\usepackage{booktabs}
\usepackage{array}
\usepackage{listings}

% Color and links
\usepackage[dvipsnames]{xcolor}
\definecolor{brandprimary}{HTML}{030391}
\definecolor{brandsecondary}{HTML}{1488D8}
\usepackage[
  colorlinks=true,
  linkcolor=brandprimary,
  citecolor=brandsecondary,
  urlcolor=brandsecondary
]{hyperref}

% Headers and footers
\usepackage{fancyhdr}
\setlength{\headheight}{36pt}
\pagestyle{fancy}
\fancyhf{}
\fancyhead[L]{%
  \parbox[c][0.42in][c]{0.42\textwidth}{%
    \raggedright\small\textbf{\coursecode}\\\coursename}}
\fancyhead[R]{%
  \parbox[c][0.42in][c]{0.42\textwidth}{%
    \raggedleft\includegraphics[height=0.42in]{assets/hcmut_logo.png}}}
\fancyfoot[L]{\small Semester \semester}
\fancyfoot[C]{\small\nouppercase{\leftmark}}
\fancyfoot[R]{\small Page \thepage}
\renewcommand{\sectionmark}[1]{\markboth{#1}{}}

% -----------------------------------------------------------------------------
% Course information -- change these values once per course
% -----------------------------------------------------------------------------
\newcommand{\universityname}{HCMC University of Technology}
\newcommand{\faculty}{Faculty of Computer Science and Engineering}
\newcommand{\coursecode}{CO5059}
\newcommand{\coursename}{Advanced System Architectures}
\newcommand{\semester}{261}
\newcommand{\instructor}{Tran Ngoc Thinh}
\newcommand{\instructoremail}{}% leave empty to hide
\newcommand{\notetaker}{Phan Nhat Minh}
\newcommand{\notetakeremail}{minh.phannhat95@hcmut.edu.vn}
\newcommand{\coursedescription}{Quantitative design and evaluation of computer
systems: performance and cost, instruction set architectures, pipelining and
its hazards, memory hierarchy, multiprocessors, superscalar/VLIW processors,
and AI \& security chips.}
\newcommand{\emaillink}[1]{%
  \href{mailto:#1}{\textcolor{brandprimary}{\texttt{#1}}}}

% Colored definitions, results, examples, proofs, and notes.
\input{\notesdir/colored-boxes}

% -----------------------------------------------------------------------------
% Course-specific notation
% -----------------------------------------------------------------------------
\newcommand{\CPI}{\mathrm{CPI}}

\lstdefinestyle{mips}{
  basicstyle=\ttfamily\small,
  keywordstyle=\bfseries\color{brandprimary},
  morekeywords={lw,sw,add,sub,and,or,slt,beq,bne,beqz,xor,nop,addi,addiu,addu,
    subu,li,la,jal,jr,jalr,sll,srl,sra,sllv,srlv,srav,andi,ori,xori,nor,slti,
    sltiu,sltu,lui,mult,multu,div,divu,mfhi,mflo,ll,sc,syscall,move,lh,lhu,lb,
    lbu,sh,sb,blez,bgtz,bltz,bgez,j,mul},
  showstringspaces=false,
  columns=fullflexible,
  keepspaces=true,
}
\lstset{style=mips}

% Creates a consistently formatted lecture heading.
% Usage: \lecture{number}{title}{date}
\newcommand{\lecture}[3]{%
  \section{Lecture #1: #2}
  \label{lec:#1}
  \begin{center}
    \small\textit{#3}
  \end{center}
}


% Cover page + contents page.  Usage: \makecover{Big title}{Description}
\newcommand{\makecover}[2]{%
\pagenumbering{roman}
\begin{titlepage}
  \thispagestyle{empty}

  \begin{center}
    \includegraphics[width=0.82\textwidth]{assets/hcmut_banner_en.png}

    \vspace{0.35cm}
    {\arialfont\large\bfseries\color{brandprimary}\faculty\par}
  \end{center}

  \vfill
  \begin{center}
    {\Huge\bfseries #1\par}
    \vspace{0.55cm}
    {\LARGE\bfseries\color{brandprimary}\coursename\par}
    \vspace{0.2cm}
    {\arialfont\large\bfseries\color{brandsecondary}\coursecode\par}

    \vspace{0.65cm}
    {\color{brandsecondary}\rule{0.38\textwidth}{1.2pt}}
    \vspace{0.65cm}

    {\large\textbf{Instructor:} \instructor
      \ifx\instructoremail\empty\else\ {\small(\emaillink{\instructoremail})}\fi\par}
    \vspace{0.3cm}
    {\large\textbf{Notes by:} \notetaker\
      {\small(\emaillink{\notetakeremail})}\par}
    \vspace{0.3cm}
    {\large\textbf{Semester:} \semester\par}

    \vspace{0.8cm}
    \begin{minipage}{0.72\textwidth}
      \centering\itshape\color{black!70}
      #2
    \end{minipage}
  \end{center}

  \vfill
  \begin{center}
    \small\color{black!60}Last updated: \today
  \end{center}
\end{titlepage}

% -----------------------------------------------------------------------------
% Table of contents -- kept on its own page
% -----------------------------------------------------------------------------
\clearpage
\setcounter{page}{2}
\thispagestyle{empty}
\hypersetup{linkcolor=black}
\setcounter{tocdepth}{2}
\tableofcontents
\hypersetup{linkcolor=brandprimary}
\clearpage
\pagenumbering{arabic}
}


\begin{document}
\sloppy
\makecover{Quiz Review: Performance}{Lecture 1 -- CPU time, CPI, speedup,
Amdahl's Law, MIPS/MFLOPS, chip cost and summarizing performance. A cheat
sheet, every calculation exercise from the slides with full solutions, the
in-class exam question solved, new exam-style questions, and extra practice
problems -- with solutions at the end.}

% ============================================================================
\section{Cheat Sheet}

\subsection{Performance and speedup}
\begin{center}
\renewcommand{\arraystretch}{2.3}
\begin{tabular}{@{}p{4.2cm}p{10.3cm}@{}}
\toprule
Quantity & Formula \\
\midrule
Performance & $\text{Perf} = \dfrac{1}{\text{ExTime}}$ \\
Speedup (A over B) & $\dfrac{\text{Perf}_A}{\text{Perf}_B} = \dfrac{\text{ExTime}_B}{\text{ExTime}_A}$ \quad (old / new) \\
``X is $n\%$ faster than Y'' & $n = 100\cdot\dfrac{\text{ExTime}_Y-\text{ExTime}_X}{\text{ExTime}_X}$ \quad (divide by the \emph{faster} time) \\
CPU time & $\text{User CPU} + \text{System CPU}$; \quad \%CPU $= \dfrac{u+s}{\text{elapsed}}$ \\
Throughput & tasks per unit time (e.g.\ passengers $\times$ mph) \\
\bottomrule
\end{tabular}
\end{center}

\subsection{The CPU equation}
\begin{theorem}[CPU time]
\[
  T = I \times \CPI \times C = \frac{I \times \CPI}{f},
  \qquad C = \frac{1}{f}
\]
\[
  \text{CPU cycles} = \sum_i \CPI_i \, C_i, \qquad
  \CPI = \frac{\text{CPU cycles}}{I} = \sum_i \CPI_i \, F_i, \qquad
  F_i = \frac{C_i}{I}
\]
\[
  \text{Fraction of time in class } i = \frac{\CPI_i\,F_i}{\CPI}
\]
\end{theorem}

\begin{center}
\begin{tabular}{@{}lccc@{}}
\toprule
Changes \dots & $I$ & $\CPI$ & $C$ \\
\midrule
Program / algorithm & X & X & \\
Compiler            & X & X & \\
ISA                 & X & X & \\
Organization        &   & X & X \\
Technology          &   &   & X \\
\bottomrule
\end{tabular}
\end{center}

\begin{note}
Unit conversions: $200\text{ MHz} \Rightarrow C = 5\text{ ns}$;
$1\text{ GHz} \Rightarrow 1\text{ ns}$; $2.5\text{ GHz} \Rightarrow
0.4\text{ ns}$. Always write $T = I\cdot\CPI / f$ with $f$ in Hz to avoid
mistakes.
\end{note}

\subsection{Amdahl's Law}
\begin{theorem}[Amdahl's Law]
If a fraction $F$ of the \emph{execution time} is sped up by $S$:
\[
  \text{ExTime}_{new} = \text{ExTime}_{old}\left[(1-F)+\frac{F}{S}\right],
  \qquad
  \text{Speedup} = \frac{1}{(1-F)+F/S} \;\le\; \frac{1}{1-F}.
\]
Several disjoint enhancements:
$\displaystyle \text{Speedup} = \frac{1}{(1-\sum F_i)+\sum F_i/S_i}$.
\end{theorem}

\begin{prop}[Required enhancement for a target speedup $T$]
\[
  \frac{1}{T} = (1-F) + \frac{F}{S}
  \quad\Longrightarrow\quad
  S = \frac{F}{\frac{1}{T}-(1-F)} .
\]
If $\frac{1}{T} \le 1-F$ (i.e.\ $T \ge \frac{1}{1-F}$), the target is
\textbf{impossible}.
\end{prop}

\begin{remark}
\textbf{Classic trap:} $F$ is a fraction of \emph{time}, not of instruction
count. If you are given an instruction mix, convert first:
$F = \CPI_i F_i / \CPI$. Alternatively, skip Amdahl and compare old vs.\ new
CPI directly with the CPU equation -- both give the same answer.
\end{remark}

\subsection{MIPS and MFLOPS}
\[
  \text{MIPS} = \frac{I}{T\times 10^6} = \frac{f}{\CPI\times10^6}
  \qquad
  \text{MFLOPS} = \frac{\#\text{FP ops}}{T\times10^6}
\]
MIPS only compares fairly with the \textbf{same program, same ISA, same
compiler}. It can rank machines \emph{opposite} to execution time.

\subsection{Cost of an integrated circuit}
\begin{align*}
  \text{Dies per wafer} &= \frac{\pi (d/2)^2}{A} - \frac{\pi d}{\sqrt{2A}}
  \qquad (d=\text{wafer diameter},\ A=\text{die area})\\
  \text{Die yield} &= \text{Wafer yield}\times\left(1+\frac{D\cdot A}{\alpha}\right)^{-\alpha},
  \qquad \alpha \approx 4\\
  \text{Die cost} &= \frac{\text{Wafer cost}}{\text{Dies per wafer}\times\text{Die yield}}\\
  \text{IC cost} &= \frac{\text{Die} + \text{Die test} + \text{Packaging} + \text{Final test}}{\text{Final test yield}}
\end{align*}
Learning curve: every doubling of volume cuts cost by $\approx10\%$.

\subsection{Summarizing several programs}
\begin{center}
\begin{tabular}{@{}lll@{}}
\toprule
Arithmetic mean & $\frac1n\sum T_i$ & tracks total time \\
Weighted mean & $\sum W_i T_i$ & result depends on the weights chosen \\
Geometric mean & $\left(\prod \frac{T_i}{T_{ref,i}}\right)^{1/n}$ & same ranking for any reference machine \\
Harmonic mean & $\frac{n}{\sum 1/R_i}$ & for averaging \emph{rates} \\
\bottomrule
\end{tabular}
\end{center}

% ============================================================================
\section{Exercises from the Slides (with Solutions)}

\begin{exercise}[Unix \texttt{time}]
The command prints \texttt{90.7u 12.9s 2:39 65\%}. Interpret it.
\end{exercise}
\begin{solution}
User CPU $=90.7$~s, system CPU $=12.9$~s, elapsed $=2{:}39=159$~s. CPU
time $=103.6$~s, and $103.6/159 = 65\%$ of the elapsed time. The other
35\% is I/O or time spent running other programs.
\end{solution}

\begin{exercise}[Two notions of performance]
Boeing 747: 6.5~h DC--Paris, 610~mph, 470 passengers. Concorde: 3~h,
1350~mph, 132 passengers. Compare latency and throughput.
\end{exercise}
\begin{solution}
Latency: Concorde is $6.5/3 = 1350/610 = 2.2\times$ faster.
Throughput (passenger-miles per hour): Boeing $470\times610=286{,}700$,
Concorde $132\times1350=178{,}200$. Boeing is $286{,}700/178{,}200 =
1.6\times$ higher throughput. Which is ``faster'' depends on the metric.
\end{solution}

\begin{exercise}[Price-performance: train vs.\ plane]
Move 2400 people. Train: 4~h trip, 2400 seats, 300~Rs/person. Plane: 45~min
trip, 300 seats, 3000~Rs/person. Compare time per person and cost-performance.
\end{exercise}
\begin{solution}
Plane needs $2400/300=8$ trips $=6$~h; train needs 4~h. Time per person:
train $4\text{ h}/2400 = 6$~s, plane $6\text{ h}/2400=9$~s, so the train
has 50\% better throughput even though each plane trip is $5\times$
faster. Cost $\times$ time: train $300\times6=1800$, plane
$3000\times9=27{,}000$ Rs$\cdot$s/person: the train is $15\times$ better.
\end{solution}

\begin{exercise}[Speedup]
A program takes 1~s on machine A and 10~s on machine B. Speedup of A over B?
\end{exercise}
\begin{solution}
$\text{Speedup} = \text{ExTime}_B/\text{ExTime}_A = 10/1 = 10$: A is
10 times faster.
\end{solution}

\begin{exercise}[CPU time]
$I = 10{,}000{,}000$ instructions, $\CPI=2.5$, clock 200~MHz. CPU time?
\end{exercise}
\begin{solution}
$C = 1/200\text{ MHz} = 5\times10^{-9}$~s, so
$T = 10^7\times2.5\times5\times10^{-9} = 0.125$~s.
\end{solution}

\begin{exercise}[New compiler and faster CPU]
Same program. A new compiler gives $I'=9{,}500{,}000$ and $\CPI'=3.0$; the
new CPU runs at 300~MHz. Speedup?
\end{exercise}
\begin{solution}
$T' = 9.5\times10^6\times3.0/(300\times10^6) = 0.095$~s.
Speedup $=0.125/0.095 = 1.32$, i.e.\ 32\% faster.
\end{solution}

\begin{exercise}[Instruction classes]
$\CPI_A=1$, $\CPI_B=2$, $\CPI_C=3$. Sequence 1 has counts $(2,1,2)$,
sequence 2 has $(4,1,1)$. Find each CPI. Which is faster?
\end{exercise}
\begin{solution}
Seq.~1: cycles $=2+2+6=10$, $I=5$, $\CPI=2.0$.
Seq.~2: cycles $=4+2+3=9$, $I=6$, $\CPI=1.5$.
Sequence 2 is faster (9 vs.\ 10 cycles) even though it executes \emph{more}
instructions.
\end{solution}

\begin{exercise}[RISC instruction mix]
ALU 50\% (CPI 1), Load 20\% (CPI 5), Store 10\% (CPI 3), Branch 20\% (CPI 2).
Find the CPI and the fraction of time spent in each class.
\end{exercise}
\begin{solution}
$\CPI = 0.5+1.0+0.3+0.4 = 2.2$. Time fractions: ALU $0.5/2.2=23\%$, Load
$1.0/2.2=45\%$, Store $0.3/2.2=14\%$, Branch $0.4/2.2=18\%$.
\end{solution}

\begin{exercise}[$n\%$ faster]
Y takes 15~s, X takes 10~s. How much faster is X?
\end{exercise}
\begin{solution}
$n = 100\,(15-10)/10 = 50\%$.
\end{solution}

\begin{exercise}[Amdahl: floating point]
FP instructions become $2\times$ faster, but only 10\% of the time is FP.
Overall speedup?
\end{exercise}
\begin{solution}
$\text{ExTime}_{new} = (0.9 + 0.1/2)\,\text{ExTime}_{old} =
0.95\,\text{ExTime}_{old}$, so Speedup $=1/0.95 = 1.053$.
\end{solution}

\begin{exercise}[Amdahl: load CPI]
In the RISC mix above, load CPI improves from 5 to 2. Speedup?
\end{exercise}
\begin{solution}
Loads take $F=45\%$ of the time (not 20\%), $S=5/2=2.5$:
$\text{Speedup}=1/(0.55+0.45/2.5) = 1/0.73 = 1.37$.
Check with CPI: new $\CPI = 0.5+0.2\times2+0.3+0.4 = 1.6$, and
$2.2/1.6 = 1.37$.
\end{solution}

\begin{exercise}[How fast must multiply be?]
A program runs in 100~s, of which multiply takes 80~s. How much faster must
multiply be to make the program (a) $4\times$ faster, (b) $5\times$ faster?
\end{exercise}
\begin{solution}
(a) Target 25~s: $25 = 20 + 80/n \Rightarrow n = 16$.
(b) Target 20~s: $20 = 20 + 80/n \Rightarrow 80/n = 0$: impossible. The
unaffected 20~s already equals the target ($1/(1-F) = 1/0.2 = 5$ is the
limit and cannot be reached).
\end{solution}

\begin{exercise}[Multiple enhancements]
$S_1=10, F_1=20\%$; $S_2=15, F_2=15\%$; $S_3=30, F_3=10\%$ (disjoint
parts). Overall speedup?
\end{exercise}
\begin{solution}
$\text{Speedup} = 1/\big(0.55 + 0.02 + 0.01 + 0.0033\big) = 1/0.5833 =
1.71$.
\end{solution}

\begin{exercise}[MIPS can mislead]
Clock 100~MHz; CPI of classes A, B, C $=1,2,3$. Compiler 1 produces
$(5,1,1)$ million instructions; compiler 2 produces $(10,1,1)$ million.
Compare MIPS and CPU time.
\end{exercise}
\begin{solution}
Compiler 1: $\CPI = 10/7 = 1.43$, MIPS $=100/1.43 = 70$, time
$=7\times10^6\times1.43/10^8 = 0.10$~s.
Compiler 2: $\CPI = 15/12 = 1.25$, MIPS $= 80$, time
$=12\times10^6\times1.25/10^8 = 0.15$~s.
Compiler 2 has the higher MIPS but is \emph{slower}.
\end{solution}

\begin{exercise}[Dies per wafer]
Die $0.7$~cm on a side, wafer diameter 30~cm. How many dies?
\end{exercise}
\begin{solution}
$A=0.49\text{ cm}^2$:
$\pi(15)^2/0.49 - \pi\cdot30/\sqrt{0.98} = 1442.6 - 95.2 \approx 1347$ dies.
\end{solution}

\begin{exercise}[Die yield]
Same die, defect density $0.6$/cm$^2$, $\alpha=4$, wafer yield 100\%.
\end{exercise}
\begin{solution}
$(1 + 0.6\times0.49/4)^{-4} = 1.0735^{-4} = 0.75$.
\end{solution}

\begin{exercise}[Weighted means]
P1/P2 take 1/1000~s on A, 10/100~s on B, 20/20~s on C. Rank the machines
under weights $(0.5,0.5)$, $(0.909,0.091)$, $(0.999,0.001)$.
\end{exercise}
\begin{solution}
\begin{tabular}[t]{@{}lccc@{}}
Weights & A & B & C \\
\hline
$.5/.5$     & 500.5 & 55.0  & \textbf{20} \\
$.909/.091$ & 91.91 & \textbf{18.19} & 20 \\
$.999/.001$ & \textbf{2.0} & 10.09 & 20 \\
\end{tabular}

Each weighting picks a different winner (C, then B, then A). (The slide
prints 91.82 and 18.8 for the middle row; the exact values are 91.91 and
18.19 -- the ranking is the same.)
\end{solution}

% ============================================================================
\section{Exam-Style Questions}

The first question was done in class; its solution follows it directly. The
others are new questions in the same style -- their solutions are in the last
section.

\begin{exercise}[In class -- two implementations of one ISA]\label{x:class}
An instruction set has four classes A, B, C, D. Implementation M1 runs at
2.4~GHz and M2 at 3.0~GHz. Cycles per instruction and class frequencies in a
typical program:
\begin{center}
\begin{tabular}{@{}cccc@{}}
\toprule
Class & M1 cycles/instr. & M2 cycles/instr. & Frequency \\
\midrule
A & 2 & 3 & 20\% \\
B & 8 & 4 & 50\% \\
C & 1 & 10 & 10\% \\
D & 4 & 2 & 20\% \\
\bottomrule
\end{tabular}
\end{center}
\begin{enumerate}[label=(\alph*),leftmargin=*]
  \item Calculate the average CPI for M1.
  \item Calculate the MIPS rating for M1.
  \item Calculate the average CPI for M2.
  \item Calculate the MIPS rating for M2.
  \item A program executes $23.7\times10^9$ instructions (on either
  implementation). How long does it take on M1, to the nearest millisecond?
  \item Repeat (e) for M2.
\end{enumerate}
\end{exercise}
\begin{solution}
\leavevmode
\begin{enumerate}[label=(\alph*),leftmargin=*]
  \item $\CPI_{M1} = 0.2(2) + 0.5(8) + 0.1(1) + 0.2(4) = 0.4+4.0+0.1+0.8 =
  5.3$.
  \item $\text{MIPS}_{M1} = \dfrac{f}{\CPI\times10^6} =
  \dfrac{2.4\times10^9}{5.3\times10^6} = 452.8$.
  \item $\CPI_{M2} = 0.2(3) + 0.5(4) + 0.1(10) + 0.2(2) = 0.6+2.0+1.0+0.4 =
  4.0$.
  \item $\text{MIPS}_{M2} = \dfrac{3.0\times10^9}{4.0\times10^6} = 750$.
  \item $T_{M1} = \dfrac{I\times\CPI}{f} = \dfrac{23.7\times10^9\times5.3}
  {2.4\times10^9} = 52.3375\text{ s} \approx 52{,}338\text{ ms}$.
  \item $T_{M2} = \dfrac{23.7\times10^9\times4.0}{3.0\times10^9} = 31.6\text{
  s} = 31{,}600\text{ ms}$.
\end{enumerate}
M2 is $52.3375/31.6 = 1.66\times$ faster. Here the MIPS ratio
($750/452.8 = 1.66$) agrees with the time ratio only because both machines
run the \emph{same} instruction count (same ISA, same program).
\end{solution}

\begin{exercise}[Two processors, one program]\label{x:p1p2}
An ISA has three classes. Processor P1 runs at 2.0~GHz and P2 at 2.5~GHz:
\begin{center}
\begin{tabular}{@{}cccc@{}}
\toprule
Class & P1 CPI & P2 CPI & Frequency \\
\midrule
A & 1 & 2 & 50\% \\
B & 3 & 2 & 30\% \\
C & 5 & 6 & 20\% \\
\bottomrule
\end{tabular}
\end{center}
\begin{enumerate}[label=(\alph*),leftmargin=*]
  \item Average CPI of P1.
  \item MIPS rating of P1.
  \item Average CPI of P2.
  \item MIPS rating of P2.
  \item A program executes $15\times10^9$ instructions. Find its execution
  time on each processor. Which is faster, and by what percentage?
  \item P1's designers reduce the CPI of class C from 5 to 3. Find P1's new
  execution time and check the speedup with Amdahl's Law. Which processor is
  faster now?
\end{enumerate}
\end{exercise}

\begin{exercise}[Two compilers]\label{x:compilers}
A 3~GHz processor has classes A, B, C with CPI 1, 2 and 4. Two compilers
produce these instruction counts (in billions) for the same program:
\begin{center}
\begin{tabular}{@{}lccc@{}}
\toprule
Compiler & A & B & C \\
\midrule
X & 6 & 2 & 2 \\
Y & 12 & 2 & 1.5 \\
\bottomrule
\end{tabular}
\end{center}
\begin{enumerate}[label=(\alph*),leftmargin=*]
  \item Total instruction count and average CPI for each compiler.
  \item MIPS rating for each compiler.
  \item CPU time for each compiler.
  \item Which compiler would a buyer choose by MIPS, and which is really
  faster? By what percentage?
  \item A hardware change halves the CPI of class C (4 $\to$ 2). For
  compiler X's code, find the new CPU time.
  \item Verify (e) with Amdahl's Law.
\end{enumerate}
\end{exercise}

\begin{exercise}[Cost of a die]\label{x:die}
A 30~cm wafer costs \$6000. The die area is $1.2\text{ cm}^2$, the defect
density $0.5/\text{cm}^2$, $\alpha = 4$, and the wafer yield is 100\%.
\begin{enumerate}[label=(\alph*),leftmargin=*]
  \item How many dies fit on the wafer?
  \item What is the die yield?
  \item How many good dies does a wafer produce?
  \item What is the cost per good die?
  \item A new process shrinks the die to $0.8\text{ cm}^2$ but the wafer costs
  \$7000. Repeat (a)--(d).
  \item Explain why the cost per die falls much more than the area.
\end{enumerate}
\end{exercise}

% ============================================================================
\section{Practice Problems}

Try these before looking at the solutions in the next section.

\begin{exercise}\label{p:time}
A program executes $2\times10^9$ instructions with $\CPI=1.5$ on a
2.5~GHz processor. What is its CPU time?
\end{exercise}

\begin{exercise}\label{p:compare}
The same program runs on processor B with $\CPI = 1.2$ at 2.2~GHz. Which
processor is faster, and by how much?
\end{exercise}

\begin{exercise}\label{p:clock}
A program runs in 10~s on a 2~GHz computer A. You are designing computer B to
run it in 6~s. The faster clock forces B to need $1.2\times$ as many clock
cycles as A. What clock rate must B have?
\end{exercise}

\begin{exercise}\label{p:compilers}
Classes A, B, C have CPI 1, 2, 3; clock 1~GHz. Compiler X emits (4, 2, 1)
million instructions of each class; compiler Y emits (2, 2, 2) million. For
each: instruction count, CPI, CPU time, MIPS. Which is faster?
\end{exercise}

\begin{exercise}\label{p:mix}
ALU 45\% (CPI 1), Load 25\% (CPI 4), Store 10\% (CPI 3), Branch 20\% (CPI 2).
Find the CPI and the fraction of time spent in each class.
\end{exercise}

\begin{exercise}\label{p:choose}
For the mix in Problem~\ref{p:mix}, you can either (a) reduce load CPI from 4
to 2, or (b) reduce branch CPI from 2 to 1. Which gives the larger speedup?
\end{exercise}

\begin{exercise}\label{p:vector}
40\% of a program's time is vectorizable and a vector unit runs that part
$10\times$ faster. What is the overall speedup? What is the best possible
speedup from vectorization alone?
\end{exercise}

\begin{exercise}\label{p:target}
FP code is 70\% of execution time. How much faster must the FP unit be for an
overall speedup of 2? Of 3.5?
\end{exercise}

\begin{exercise}\label{p:twoenh}
Enhancement E1 speeds up 20\% of the time by $10\times$; E2 speeds up 50\% of
the time by $1.6\times$ (different parts). Which one alone is better? What do
you get with both?
\end{exercise}

\begin{exercise}\label{p:percent}
(a) X takes 8~s and Y takes 12~s. How much faster is X?
(b) X is 25\% faster than Y, and Y takes 10~s. How long does X take?
\end{exercise}

\begin{exercise}\label{p:compilerchange}
A new compiler reduces instruction count by 20\% but raises CPI by 10\%.
Speedup?
\end{exercise}

\begin{exercise}\label{p:mips}
A 500~MHz machine runs a mix of 60\% (CPI 1), 30\% (CPI 2), 10\% (CPI 4).
Find CPI, MIPS rating, and the time for $10^9$ instructions.
\end{exercise}

\begin{exercise}\label{p:unix}
\texttt{time} prints \texttt{45.3u 4.7s 1:40 50\%}. Check the 50\% and find
what fraction of CPU time is user time.
\end{exercise}

\begin{exercise}\label{p:die}
Wafer diameter 20~cm, wafer cost \$1000, defect density 0.8/cm$^2$,
$\alpha=4$, wafer yield 100\%. Find dies per wafer, die yield and die cost
for (a) a 1~cm$^2$ die and (b) a 0.5~cm$^2$ die.
\end{exercise}

\begin{exercise}\label{p:geo}
Times (s) for P1, P2: machine A $(2, 20)$, machine B $(4, 10)$, reference R
$(4, 20)$. Compare A and B by total time and by geometric mean of
normalized times. Repeat the geometric mean using A as the reference.
\end{exercise}

% ============================================================================
\section{Solutions to the Practice Problems}

\begin{solution}[\ref{p:time}]
$T = 2\times10^9 \times 1.5 / (2.5\times10^9) = 1.2$~s.
\end{solution}

\begin{solution}[\ref{p:compare}]
$T_B = 2\times10^9\times1.2/(2.2\times10^9) = 1.091$~s. B is faster:
$1.2/1.091 = 1.10$, i.e.\ 10\% faster.
\end{solution}

\begin{solution}[\ref{p:clock}]
Cycles on A $= 10\times2\times10^9 = 2\times10^{10}$. B needs
$2.4\times10^{10}$ cycles in 6~s: $f_B = 2.4\times10^{10}/6 = 4$~GHz.
\end{solution}

\begin{solution}[\ref{p:compilers}]
X: $I=7$M, cycles $=4+4+3=11$M, $\CPI=1.57$, $T=11$~ms, MIPS
$=1000/1.57=636$.
Y: $I=6$M, cycles $=2+4+6=12$M, $\CPI=2.0$, $T=12$~ms, MIPS $=500$.
X is faster ($12/11=1.09\times$) even though it executes more instructions.
\end{solution}

\begin{solution}[\ref{p:mix}]
$\CPI = 0.45+1.00+0.30+0.40 = 2.15$. Time fractions: ALU $20.9\%$, Load
$46.5\%$, Store $14.0\%$, Branch $18.6\%$.
\end{solution}

\begin{solution}[\ref{p:choose}]
(a) New $\CPI = 2.15-0.25\times2 = 1.65$, speedup $=2.15/1.65=1.30$.
(b) New $\CPI = 2.15-0.20\times1 = 1.95$, speedup $=1.10$.
Option (a) wins: loads are the common case in \emph{time} (46.5\%).
Amdahl check for (a): $1/(0.535+0.465/2) = 1.30$.
\end{solution}

\begin{solution}[\ref{p:vector}]
$1/(0.6+0.4/10) = 1/0.64 = 1.56$. Even an infinitely fast vector unit gives
at most $1/0.6 = 1.67$.
\end{solution}

\begin{solution}[\ref{p:target}]
For 2: $0.5 = 0.3 + 0.7/S \Rightarrow S = 3.5$. For 3.5: $1/3.5=0.286 <
0.3$, impossible; the maximum is $1/0.3=3.33$.
\end{solution}

\begin{solution}[\ref{p:twoenh}]
E1: $1/(0.8+0.02) = 1.22$. E2: $1/(0.5+0.3125) = 1.23$. E2 is (slightly)
better. Both: $1/(0.3+0.02+0.3125)=1/0.6325 = 1.58$.
\end{solution}

\begin{solution}[\ref{p:percent}]
(a) $n = 100\times(12-8)/8 = 50\%$.
(b) $\text{ExTime}_Y/\text{ExTime}_X = 1.25 \Rightarrow \text{ExTime}_X
= 10/1.25 = 8$~s (not $7.5$~s).
\end{solution}

\begin{solution}[\ref{p:compilerchange}]
$T'/T = 0.8\times1.1 = 0.88$, so speedup $=1/0.88 = 1.14$.
\end{solution}

\begin{solution}[\ref{p:mips}]
$\CPI = 0.6+0.6+0.4 = 1.6$; MIPS $=500/1.6 = 312.5$;
$T = 10^9\times1.6/(5\times10^8) = 3.2$~s ($=10^9/312.5\times10^6$).
\end{solution}

\begin{solution}[\ref{p:unix}]
Elapsed $=100$~s; $(45.3+4.7)/100 = 50\%$. User share of CPU time
$=45.3/50 = 90.6\%$.
\end{solution}

\begin{solution}[\ref{p:die}]
(a) Dies $= \pi\cdot100/1 - \pi\cdot20/\sqrt2 = 314.2-44.4 = 269$.
Yield $=(1+0.8/4)^{-4} = 1.2^{-4} = 0.482$. Die cost
$=1000/(269\times0.482) = \$7.71$.
(b) Dies $=628.3 - 62.8 = 565$. Yield $=1.1^{-4}=0.683$. Die cost
$=1000/(565\times0.683) = \$2.59$.
Halving the area cuts the die cost by about $3\times$: more dies fit
\emph{and} each is less likely to contain a defect.
\end{solution}

\begin{solution}[\ref{p:geo}]
Total time: A $=22$~s, B $=14$~s, so B is better for this workload.
Normalized to R: A $=(0.5, 1.0)$, B $=(1.0, 0.5)$; both geometric means
are $\sqrt{0.5}=0.707$, a tie. Normalized to A: A $=(1,1)\to1$, B
$=(2, 0.5)\to1$, again a tie. The geometric mean is consistent across
references, but it does not track total execution time.
\end{solution}

% ============================================================================
\section{Solutions to the Exam-Style Questions}

\begin{solution}[\ref{x:p1p2}]
\leavevmode
\begin{enumerate}[label=(\alph*),leftmargin=*]
  \item $\CPI_{P1} = 0.5(1)+0.3(3)+0.2(5) = 0.5+0.9+1.0 = 2.4$.
  \item $\text{MIPS}_{P1} = 2000/2.4 = 833.3$.
  \item $\CPI_{P2} = 0.5(2)+0.3(2)+0.2(6) = 1.0+0.6+1.2 = 2.8$.
  \item $\text{MIPS}_{P2} = 2500/2.8 = 892.9$.
  \item $T_{P1} = 15\times10^9\times2.4/(2\times10^9) = 18$~s;
  $T_{P2} = 15\times10^9\times2.8/(2.5\times10^9) = 16.8$~s. P2 is faster:
  $18/16.8 = 1.071$, i.e.\ 7.1\% faster.
  \item New $\CPI_{P1} = 0.5+0.9+0.2(3) = 2.0$, $T = 15$~s, speedup $18/15 =
  1.2$. Amdahl: class C took $1.0/2.4 = 41.7\%$ of the time, $S = 5/3$:
  $1/(0.583 + 0.417\times0.6) = 1/0.833 = 1.2$. P1 is now faster than P2 by
  $16.8/15 = 1.12$ (12\%).
\end{enumerate}
\end{solution}

\begin{solution}[\ref{x:compilers}]
\leavevmode
\begin{enumerate}[label=(\alph*),leftmargin=*]
  \item X: $I = 10\times10^9$, cycles $= 6+4+8 = 18\times10^9$, $\CPI = 1.8$.
  Y: $I = 15.5\times10^9$, cycles $= 12+4+6 = 22\times10^9$, $\CPI = 22/15.5
  = 1.42$.
  \item X: $3000/1.8 = 1667$ MIPS. Y: $3000/1.42 = 2114$ MIPS.
  \item X: $18\times10^9/(3\times10^9) = 6$~s. Y: $22\times10^9/(3\times10^9)
  = 7.33$~s.
  \item MIPS favours Y, but X is faster: $7.33/6 = 1.22$, i.e.\ 22\% faster.
  Y executes many more (cheap) instructions, which raises MIPS but also the
  total time.
  \item Cycles $= 6 + 4 + 2\times2 = 14\times10^9$, $T = 4.67$~s.
  \item Class C took $8/18 = 44.4\%$ of X's time, $S = 2$:
  $1/(0.556 + 0.222) = 1.286 = 6/4.67$. \checkmark
\end{enumerate}
\end{solution}

\begin{solution}[\ref{x:die}]
\leavevmode
\begin{enumerate}[label=(\alph*),leftmargin=*]
  \item $\dfrac{\pi(15)^2}{1.2} - \dfrac{\pi\cdot30}{\sqrt{2.4}} = 589.0 -
  60.8 = 528$ dies.
  \item $(1 + 0.5\times1.2/4)^{-4} = 1.15^{-4} = 0.572$.
  \item $528\times0.572 \approx 302$ good dies.
  \item $6000/301.9 \approx \$19.9$.
  \item Dies $= 883.6 - 74.5 = 809$; yield $= (1+0.1)^{-4} = 0.683$; good dies
  $\approx 553$; cost $= 7000/552.6 \approx \$12.7$.
  \item A smaller die gives more dies per wafer (fewer wasted at the edge, too)
  \emph{and} a higher yield, because each die is less likely to contain a
  defect. The two effects multiply, so cost falls faster than area (here a
  33\% smaller die is 36\% cheaper even though the wafer costs 17\% more).
\end{enumerate}
\end{solution}

\end{document}
